% Options for packages loaded elsewhere
\PassOptionsToPackage{unicode}{hyperref}
\PassOptionsToPackage{hyphens}{url}
\documentclass[
  english,
  man]{apa6}
\usepackage{xcolor}
\usepackage{amsmath,amssymb}
\setcounter{secnumdepth}{-\maxdimen} % remove section numbering
\usepackage{iftex}
\ifPDFTeX
  \usepackage[T1]{fontenc}
  \usepackage[utf8]{inputenc}
  \usepackage{textcomp} % provide euro and other symbols
\else % if luatex or xetex
  \usepackage{unicode-math} % this also loads fontspec
  \defaultfontfeatures{Scale=MatchLowercase}
  \defaultfontfeatures[\rmfamily]{Ligatures=TeX,Scale=1}
\fi
\usepackage{lmodern}
\ifPDFTeX\else
  % xetex/luatex font selection
\fi
% Use upquote if available, for straight quotes in verbatim environments
\IfFileExists{upquote.sty}{\usepackage{upquote}}{}
\IfFileExists{microtype.sty}{% use microtype if available
  \usepackage[]{microtype}
  \UseMicrotypeSet[protrusion]{basicmath} % disable protrusion for tt fonts
}{}
\makeatletter
\@ifundefined{KOMAClassName}{% if non-KOMA class
  \IfFileExists{parskip.sty}{%
    \usepackage{parskip}
  }{% else
    \setlength{\parindent}{0pt}
    \setlength{\parskip}{6pt plus 2pt minus 1pt}}
}{% if KOMA class
  \KOMAoptions{parskip=half}}
\makeatother
% Make \paragraph and \subparagraph free-standing
\makeatletter
\ifx\paragraph\undefined\else
  \let\oldparagraph\paragraph
  \renewcommand{\paragraph}{
    \@ifstar
      \xxxParagraphStar
      \xxxParagraphNoStar
  }
  \newcommand{\xxxParagraphStar}[1]{\oldparagraph*{#1}\mbox{}}
  \newcommand{\xxxParagraphNoStar}[1]{\oldparagraph{#1}\mbox{}}
\fi
\ifx\subparagraph\undefined\else
  \let\oldsubparagraph\subparagraph
  \renewcommand{\subparagraph}{
    \@ifstar
      \xxxSubParagraphStar
      \xxxSubParagraphNoStar
  }
  \newcommand{\xxxSubParagraphStar}[1]{\oldsubparagraph*{#1}\mbox{}}
  \newcommand{\xxxSubParagraphNoStar}[1]{\oldsubparagraph{#1}\mbox{}}
\fi
\makeatother
\usepackage{graphicx}
\makeatletter
\newsavebox\pandoc@box
\newcommand*\pandocbounded[1]{% scales image to fit in text height/width
  \sbox\pandoc@box{#1}%
  \Gscale@div\@tempa{\textheight}{\dimexpr\ht\pandoc@box+\dp\pandoc@box\relax}%
  \Gscale@div\@tempb{\linewidth}{\wd\pandoc@box}%
  \ifdim\@tempb\p@<\@tempa\p@\let\@tempa\@tempb\fi% select the smaller of both
  \ifdim\@tempa\p@<\p@\scalebox{\@tempa}{\usebox\pandoc@box}%
  \else\usebox{\pandoc@box}%
  \fi%
}
% Set default figure placement to htbp
\def\fps@figure{htbp}
\makeatother
% definitions for citeproc citations
\NewDocumentCommand\citeproctext{}{}
\NewDocumentCommand\citeproc{mm}{%
  \begingroup\def\citeproctext{#2}\cite{#1}\endgroup}
\makeatletter
 % allow citations to break across lines
 \let\@cite@ofmt\@firstofone
 % avoid brackets around text for \cite:
 \def\@biblabel#1{}
 \def\@cite#1#2{{#1\if@tempswa , #2\fi}}
\makeatother
\newlength{\cslhangindent}
\setlength{\cslhangindent}{1.5em}
\newlength{\csllabelwidth}
\setlength{\csllabelwidth}{3em}
\newenvironment{CSLReferences}[2] % #1 hanging-indent, #2 entry-spacing
 {\begin{list}{}{%
  \setlength{\itemindent}{0pt}
  \setlength{\leftmargin}{0pt}
  \setlength{\parsep}{0pt}
  % turn on hanging indent if param 1 is 1
  \ifodd #1
   \setlength{\leftmargin}{\cslhangindent}
   \setlength{\itemindent}{-1\cslhangindent}
  \fi
  % set entry spacing
  \setlength{\itemsep}{#2\baselineskip}}}
 {\end{list}}
\usepackage{calc}
\newcommand{\CSLBlock}[1]{\hfill\break\parbox[t]{\linewidth}{\strut\ignorespaces#1\strut}}
\newcommand{\CSLLeftMargin}[1]{\parbox[t]{\csllabelwidth}{\strut#1\strut}}
\newcommand{\CSLRightInline}[1]{\parbox[t]{\linewidth - \csllabelwidth}{\strut#1\strut}}
\newcommand{\CSLIndent}[1]{\hspace{\cslhangindent}#1}
\ifLuaTeX
\usepackage[bidi=basic]{babel}
\else
\usepackage[bidi=default]{babel}
\fi
% get rid of language-specific shorthands (see #6817):
\let\LanguageShortHands\languageshorthands
\def\languageshorthands#1{}
\ifLuaTeX
  \usepackage[english]{selnolig} % disable illegal ligatures
\fi
\setlength{\emergencystretch}{3em} % prevent overfull lines
\providecommand{\tightlist}{%
  \setlength{\itemsep}{0pt}\setlength{\parskip}{0pt}}
% Manuscript styling
\usepackage{upgreek}
\captionsetup{font=singlespacing,justification=justified}

% Table formatting
\usepackage{longtable}
\usepackage{lscape}
% \usepackage[counterclockwise]{rotating}   % Landscape page setup for large tables
\usepackage{multirow}		% Table styling
\usepackage{tabularx}		% Control Column width
\usepackage[flushleft]{threeparttable}	% Allows for three part tables with a specified notes section
\usepackage{threeparttablex}            % Lets threeparttable work with longtable

% Create new environments so endfloat can handle them
% \newenvironment{ltable}
%   {\begin{landscape}\centering\begin{threeparttable}}
%   {\end{threeparttable}\end{landscape}}
\newenvironment{lltable}{\begin{landscape}\centering\begin{ThreePartTable}}{\end{ThreePartTable}\end{landscape}}

% Enables adjusting longtable caption width to table width
% Solution found at http://golatex.de/longtable-mit-caption-so-breit-wie-die-tabelle-t15767.html
\makeatletter
\newcommand\LastLTentrywidth{1em}
\newlength\longtablewidth
\setlength{\longtablewidth}{1in}
\newcommand{\getlongtablewidth}{\begingroup \ifcsname LT@\roman{LT@tables}\endcsname \global\longtablewidth=0pt \renewcommand{\LT@entry}[2]{\global\advance\longtablewidth by ##2\relax\gdef\LastLTentrywidth{##2}}\@nameuse{LT@\roman{LT@tables}} \fi \endgroup}

% \setlength{\parindent}{0.5in}
% \setlength{\parskip}{0pt plus 0pt minus 0pt}

% Overwrite redefinition of paragraph and subparagraph by the default LaTeX template
% See https://github.com/crsh/papaja/issues/292
\makeatletter
\renewcommand{\paragraph}{\@startsection{paragraph}{4}{\parindent}%
  {0\baselineskip \@plus 0.2ex \@minus 0.2ex}%
  {-1em}%
  {\normalfont\normalsize\bfseries\itshape\typesectitle}}

\renewcommand{\subparagraph}[1]{\@startsection{subparagraph}{5}{1em}%
  {0\baselineskip \@plus 0.2ex \@minus 0.2ex}%
  {-\z@\relax}%
  {\normalfont\normalsize\itshape\hspace{\parindent}{#1}\textit{\addperi}}{\relax}}
\makeatother

\makeatletter
\usepackage{etoolbox}
\patchcmd{\maketitle}
  {\section{\normalfont\normalsize\abstractname}}
  {\section*{\normalfont\normalsize\abstractname}}
  {}{\typeout{Failed to patch abstract.}}
\patchcmd{\maketitle}
  {\section{\protect\normalfont{\@title}}}
  {\section*{\protect\normalfont{\@title}}}
  {}{\typeout{Failed to patch title.}}
\makeatother

\usepackage{xpatch}
\makeatletter
\xapptocmd\appendix
  {\xapptocmd\section
    {\addcontentsline{toc}{section}{\appendixname\ifoneappendix\else~\theappendix\fi: #1}}
    {}{\InnerPatchFailed}%
  }
{}{\PatchFailed}
\makeatother
\keywords{keywords\newline\indent Word count: X}
\DeclareDelayedFloatFlavor{ThreePartTable}{table}
\DeclareDelayedFloatFlavor{lltable}{table}
\DeclareDelayedFloatFlavor*{longtable}{table}
\makeatletter
\renewcommand{\efloat@iwrite}[1]{\immediate\expandafter\protected@write\csname efloat@post#1\endcsname{}}
\makeatother
\usepackage{csquotes}
\usepackage{bookmark}
\IfFileExists{xurl.sty}{\usepackage{xurl}}{} % add URL line breaks if available
\urlstyle{same}
\hypersetup{
  pdftitle={PSY 8960-003 Reproducibility Project: Examining the Association Between Neuroticism and Loneliness Among College Students},
  pdfauthor={Caroline Armstrong1},
  pdflang={en-EN},
  pdfkeywords={keywords},
  hidelinks,
  pdfcreator={LaTeX via pandoc}}

\title{PSY 8960-003 Reproducibility Project: Examining the Association Between Neuroticism and Loneliness Among College Students}
\author{Caroline Armstrong\textsuperscript{1}}
\date{}


\shorttitle{Reproducibility Project}

\authornote{

The authors made the following contributions. Caroline Armstrong: Conceptualization, Writing - Original Draft Preparation, Writing - Review \& Editing.

Correspondence concerning this article should be addressed to Caroline Armstrong, 75 East River Parkway, Minneapolis, MN 55455. E-mail: \href{mailto:carolinearmstrong19@gmail.com}{\nolinkurl{carolinearmstrong19@gmail.com}}

}

\affiliation{\vspace{0.5cm}\textsuperscript{1} Department of Psychology, University of Minnesota}

\abstract{%
One or two sentences providing a \textbf{basic introduction} to the field, comprehensible to a scientist in any discipline.
Two to three sentences of \textbf{more detailed background}, comprehensible to scientists in related disciplines.
One sentence clearly stating the \textbf{general problem} being addressed by this particular study.
One sentence summarizing the main result (with the words ``\textbf{here we show}'' or their equivalent).
Two or three sentences explaining what the \textbf{main result} reveals in direct comparison to what was thought to be the case previously, or how the main result adds to previous knowledge.
One or two sentences to put the results into a more \textbf{general context}.
Two or three sentences to provide a \textbf{broader perspective}, readily comprehensible to a scientist in any discipline.
}



\begin{document}
\maketitle

{[}Intro - TO COME{]}

\section{Methods}\label{methods}

\subsection{Participants}\label{participants}

Participants will be 400 university students enrolled in a Psychology course who opt to participate in a brief online survey for course credit.

\subsection{Measures}\label{measures}

\subsubsection{Demographics}\label{demographics}

Participants will be asked to indicate their age, gender, and racial and ethnic background. The exact questions and further details are shown in Appendix A.

\subsubsection{Big Five Inventory---Neuroticism Subscale}\label{big-five-inventoryneuroticism-subscale}

The Big Five Inventory---Neuroticism subscale (John, Donahue, \& Kentle, 1991) asks participants to indicate the extent to which they agree or disagree with eight statements about their personal characteristics or thinking and feeling tendencies in day-to-day life (e.g., ``I see myself as someone who is depressed, blue;'' ``I see myself as someone who worries a lot''). The response scale is an ordinal scale ranging from 1 = ``Disagree strongly'' to 5 = ``Agree strongly,'' with the intermediate response options being 2 = ``Disagree a little,''3 = ``Neutral; no opinion,'' and 4 = ``Agree a little.'' The complete scale instructions and items are shown in Appendix B.

A composite score is computed by summing the eight items.

\subsubsection{Loneliness in Context Questionnaire for College Students}\label{loneliness-in-context-questionnaire-for-college-students}

The Loneliness in Context Questionnaire for College Students (Asher \& Weeks, 2012) asks participants to indicate the frequency with which they feel lonely in various daily or weekly contexts typical for college students. Ten items tap different contexts (e.g., ``Class is a lonely place for me;'' ``I feel sad and alone at social events''), and participants respond on an ordinal scale ranging from 1 = ``Never'' to 5 = ``Always,'' with the intermediate response options being 2 = ``Hardly ever,'' 3 = ``Sometimes,'' and 4 = ``Most of the time.''

A composite score is computed by taking the average of the ten items.

\subsection{Procedure}\label{procedure}

The online survey will be developed and administered via Qualtrics. Potential participants will be informed about the study via online communications from the university's Psychology Department and may choose to complete the study from their personal device at any time during the current academic semester.

Participants will provide informed consent by indicating that they would like to proceed to the survey questions after reading a brief written explanation (on the first page of the Qualtrics survey) describing what they will be asked to do (i.e., complete an approximately five-minute survey asking them about their personality and mental state) and the possible risks (i.e., experiencing various forms of psychological discomfort while reflecting on oneself and one's feelings) and benefits (i.e., possible improvements in scientific knowledge) of their participation.

Of note, the consent will include a clause indicating that participants' de-identified data (i.e., their responses to all survey questions, including demographic information, but NOT including any information such as IP addresses that could be used to identify them) may be used by other researchers for any purpose and shared on an online, publicly available platform.

Participants who affirm their consent will proceed to the survey questions. All participants will complete the Demographic questions first. To control for possible order effects, participants will be randomized to complete either the Big Five Inventory---Neuroticism subscale or the Loneliness in Context Survey second. The third section will consist of whichever measure participants did not complete second.

Following the survey questions, participants will view a short debriefing description indicating the study aim and providing information about campus mental health resources that may be of aid if participants found that reflecting on their mental state and day-to-day loneliness caused distress or prompted them to think about seeking psychological help.

\subsection{Data analysis}\label{data-analysis}

We used R (Version 4.5.2; R Core Team, 2025) and the R-packages \emph{apaTables} (Version 2.0.8; Stanley, 2021), \emph{dplyr} (Version 1.1.4; Wickham, François, Henry, Müller, \& Vaughan, 2023), \emph{faux} (Version 1.2.3; DeBruine, 2025), \emph{ggplot2} (Version 4.0.1; Wickham, 2016), \emph{groundhog} (Version 3.2.3; Simonsohn \& Gruson, 2025), \emph{gtsummary} (Version 2.5.0; Sjoberg, Whiting, Curry, Lavery, \& Larmarange, 2021), \emph{papaja} (Version 0.1.4; Aust \& Barth, 2025), \emph{psych} (Version 2.5.6; William Revelle, 2025), \emph{readr} (Version 2.1.5; Wickham, Hester, \& Bryan, 2024), and \emph{tinylabels} (Version 0.2.5; Barth, 2025) for all our analyses.

\section{Results}\label{results}

\section{Discussion}\label{discussion}

\newpage

\section{References}\label{references}

\phantomsection\label{refs}
\begin{CSLReferences}{1}{0}
\bibitem[\citeproctext]{ref-R-papaja}
Aust, F., \& Barth, M. (2025). \emph{{papaja}: {Prepare} reproducible {APA} journal articles with {R Markdown}}. \url{https://doi.org/10.32614/CRAN.package.papaja}

\bibitem[\citeproctext]{ref-R-tinylabels}
Barth, M. (2025). \emph{{tinylabels}: Lightweight variable labels}. \url{https://doi.org/10.32614/CRAN.package.tinylabels}

\bibitem[\citeproctext]{ref-R-faux}
DeBruine, L. (2025). \emph{Faux: Simulation for factorial designs}. Zenodo. \url{https://doi.org/10.5281/zenodo.2669586}

\bibitem[\citeproctext]{ref-john1991}
John, O. P., Donahue, E. M., \& Kentle, R. L. (1991). \emph{Big Five Inventory}. {[}Database record{]}. APA PsycTests. \url{https://doi.org/10.1037/t07550-000}

\bibitem[\citeproctext]{ref-R-base}
R Core Team. (2025). \emph{R: A language and environment for statistical computing}. Vienna, Austria: R Foundation for Statistical Computing. Retrieved from \url{https://www.R-project.org/}

\bibitem[\citeproctext]{ref-R-groundhog}
Simonsohn, U., \& Gruson, H. (2025). \emph{Groundhog: Version-control for CRAN, GitHub, and GitLab packages}. \url{https://doi.org/10.32614/CRAN.package.groundhog}

\bibitem[\citeproctext]{ref-R-gtsummary}
Sjoberg, D. D., Whiting, K., Curry, M., Lavery, J. A., \& Larmarange, J. (2021). Reproducible summary tables with the gtsummary package. \emph{{The R Journal}}, \emph{13}, 570--580. \url{https://doi.org/10.32614/RJ-2021-053}

\bibitem[\citeproctext]{ref-R-apaTables}
Stanley, D. (2021). \emph{apaTables: Create american psychological association (APA) style tables}. \url{https://doi.org/10.32614/CRAN.package.apaTables}

\bibitem[\citeproctext]{ref-R-ggplot2}
Wickham, H. (2016). \emph{ggplot2: Elegant graphics for data analysis}. Springer-Verlag New York. Retrieved from \url{https://ggplot2.tidyverse.org}

\bibitem[\citeproctext]{ref-R-dplyr}
Wickham, H., François, R., Henry, L., Müller, K., \& Vaughan, D. (2023). \emph{Dplyr: A grammar of data manipulation}. \url{https://doi.org/10.32614/CRAN.package.dplyr}

\bibitem[\citeproctext]{ref-R-readr}
Wickham, H., Hester, J., \& Bryan, J. (2024). \emph{Readr: Read rectangular text data}. \url{https://doi.org/10.32614/CRAN.package.readr}

\bibitem[\citeproctext]{ref-R-psych}
William Revelle. (2025). \emph{Psych: Procedures for psychological, psychometric, and personality research}. Evanston, Illinois: Northwestern University. Retrieved from \url{https://CRAN.R-project.org/package=psych}

\end{CSLReferences}


\end{document}
